\input{../tex-inputs/latex-leseansicht-vorspann}

\section[Arthur Schnitzler, Felix Salten und Ottilie Metzeles an Gustav Schwarzkopf, 25. 5. 1900]{L04450 Arthur Schnitzler, Felix Salten und Ottilie Metzeles an Gustav
               Schwarzkopf, 25. 5. 1900}
\nopagebreak\mylabel{L04450v}
\rehead{ }\normalsize\beginnumbering\briefempfaengerindex{Schwarzkopf, Gustav@\textsc{Schwarzkopf, Gustav}!zzzSalten, Ottilie@\emph{von Ottilie Salten}!1900-05-251@{25. 5. 1900}|(be}\briefempfaengerindex{Schwarzkopf, Gustav@\textsc{Schwarzkopf, Gustav}!zzzSalten, Felix@\emph{von Felix Salten}!1900-05-251@{25. 5. 1900}|(be}\briefempfaengerindex{Schwarzkopf, Gustav@\textsc{Schwarzkopf, Gustav}!zzzSchnitzler, Arthur@\emph{von Arthur Schnitzler}!1900-05-251@{25. 5. 1900}|(be}
\toendnotes[C]{\smallbreak\pagebreak[2]}
\correspDesc{Versand  durch Arthur Schnitzler, Felix Salten, Ottilie Metzeles am 25. 5. 1900 in Puchberg am Schneeberg
\newline{}Übermittlung  am 26. 5. 1900 in Puchberg am Schneeberg
\newline{}Erhalt  durch Gustav Schwarzkopf am 27. 5. 1900 in Wien}\toendnotes[C]{\smallbreak}
\Standort{DLA, A:Schnitzler, HS.1985.1.1897.}
\physDesc{Bildpostkarte, 228 Zeichen
\newline{}Handschrift Arthur Schnitzler: Bleistift, deutsche Kurrent
\newline{}Handschrift Felix Salten: Bleistift, lateinische Kurrent
\newline{}Handschrift Ottilie Salten: Bleistift, lateinische Kurrent
\newline{}Versand: 1) Stempel: »\nobreak{}\oindex{Puchberg am Schneeberg@\textbf{Puchberg am Schneeberg}|pwk}Puchberg {[}am{]}
                                       Schneeberge, \textcolor{gray}{26}/5 00\nobreak{}«.   2) Stempel: »\nobreak{}\oindex{I., Innere Stadt@\textbf{I., Innere Stadt}|pwk}Wien \textcolor{gray}{1/1 1}, \textcolor{gray}{27}/{[}7{]} 00, 9–10\textcolor{gray}{½ V}, bestellt\nobreak{}«. }\toendnotes[C]{\smallbreak}\pstart{}{\pb}\textsc{Herrn Gustav
                     Schwarzkopf}\pend{}\pstart{}Wien\oindex{Wien@\textbf{Wien}|pw}\pend{}\pstart{}\textsc{I. Tiefer
                        Graben 23}\oindex{Wien@\textbf{Wien}!I., Innere Stadt@\textbf{I., Innere Stadt}!Tiefer Graben 23@\textbf{Tiefer Graben 23}|pw}\pend{}{\bigskip}
\pstart
           \noindent{}\centering{}{\pb}\textcolor{gray}{\textbf{Gruss vom Schneeberg\oindex{Hochschneeberg@\textbf{Hochschneeberg}|pw}.}}\pend
           \vspace{1em}
\pstart
           {\pb}25. 5. 900.\pend
           \vspace{0.5em}
\pstart
           Nicht als ob wir oben geweſen wären; wir ſind in Puchberg\oindex{Puchberg am Schneeberg@\textbf{Puchberg am Schneeberg}|pw}, wo es wirklich hübsch iſt, wenn es in dieſer Stunde auch
               regnet.\pend
           \pstart Herzlichſt Ihr \spacefill\mbox{ArthurSch}\pend{}\selectlanguage{ngerman}\vspace{1em}
\pstart
           \noindent{}\label{K_L04450-1v}\edtext{{[}hs. F. Salten:{]} Viele Grüße\spacefill\mbox{F. Salten}{ }\spacefill\mbox{{[}hs. O. Salten:{]} Otti}}{\lemma{\textnormal{\emph{Viele … Otti}}}\Cendnote{\textnormal{Die Grüße von Felix Salten\pwindex{Salten, Felix 6.\,9.\,1869 Budapest – 8.\,10.\,1945 Zürich@\textsc{Salten, Felix} (6.\,9.\,1869 Budapest – 8.\,10.\,1945 Zürich)|pwk} und Ottilie Metzeles\pwindex{Salten, Ottilie 7.\,3.\,1868 Prag – 22.\,6.\,1942 Zürich@\textsc{Salten, Ottilie} (7.\,3.\,1868 Prag – 22.\,6.\,1942 Zürich)|pwk} befinden sich entlang des linken Kartenrandes.}}}\label{K_L04450-1}\pend
           \selectlanguage{ngerman}\endnumbering\briefempfaengerindex{Schwarzkopf, Gustav@\textsc{Schwarzkopf, Gustav}!zzzSalten, Ottilie@\emph{von Ottilie Salten}!1900-05-251@{25. 5. 1900}|)be}\briefempfaengerindex{Schwarzkopf, Gustav@\textsc{Schwarzkopf, Gustav}!zzzSalten, Felix@\emph{von Felix Salten}!1900-05-251@{25. 5. 1900}|)be}\briefempfaengerindex{Schwarzkopf, Gustav@\textsc{Schwarzkopf, Gustav}!zzzSchnitzler, Arthur@\emph{von Arthur Schnitzler}!1900-05-251@{25. 5. 1900}|)be}\mylabel{L04450h}
\begin{anhang}
\end{anhang}\newcommand{\dateiname}{L04450}\newcommand{\titel}{Arthur Schnitzler, Felix Salten und Ottilie Metzeles an Gustav Schwarzkopf, 25. 5. 1900}\newcommand{\editorInnen}{Herausgegeben von Jahnke, SelmaMüller, Martin Anton}\input{../tex-inputs/latex-leseansicht-abspann}
